%% ODER: format ==         = "\mathrel{==}"
%% ODER: format /=         = "\neq "
%
%
\makeatletter
\@ifundefined{lhs2tex.lhs2tex.sty.read}%
  {\@namedef{lhs2tex.lhs2tex.sty.read}{}%
   \newcommand\SkipToFmtEnd{}%
   \newcommand\EndFmtInput{}%
   \long\def\SkipToFmtEnd#1\EndFmtInput{}%
  }\SkipToFmtEnd

\newcommand\ReadOnlyOnce[1]{\@ifundefined{#1}{\@namedef{#1}{}}\SkipToFmtEnd}
\usepackage{amstext}
\usepackage{amssymb}
\usepackage{stmaryrd}
\DeclareFontFamily{OT1}{cmtex}{}
\DeclareFontShape{OT1}{cmtex}{m}{n}
  {<5><6><7><8>cmtex8
   <9>cmtex9
   <10><10.95><12><14.4><17.28><20.74><24.88>cmtex10}{}
\DeclareFontShape{OT1}{cmtex}{m}{it}
  {<-> ssub * cmtt/m/it}{}
\newcommand{\texfamily}{\fontfamily{cmtex}\selectfont}
\DeclareFontShape{OT1}{cmtt}{bx}{n}
  {<5><6><7><8>cmtt8
   <9>cmbtt9
   <10><10.95><12><14.4><17.28><20.74><24.88>cmbtt10}{}
\DeclareFontShape{OT1}{cmtex}{bx}{n}
  {<-> ssub * cmtt/bx/n}{}
\newcommand{\tex}[1]{\text{\texfamily#1}}	% NEU

\newcommand{\Sp}{\hskip.33334em\relax}


\newcommand{\Conid}[1]{\mathit{#1}}
\newcommand{\Varid}[1]{\mathit{#1}}
\newcommand{\anonymous}{\kern0.06em \vbox{\hrule\@width.5em}}
\newcommand{\plus}{\mathbin{+\!\!\!+}}
\newcommand{\bind}{\mathbin{>\!\!\!>\mkern-6.7mu=}}
\newcommand{\rbind}{\mathbin{=\mkern-6.7mu<\!\!\!<}}% suggested by Neil Mitchell
\newcommand{\sequ}{\mathbin{>\!\!\!>}}
\renewcommand{\leq}{\leqslant}
\renewcommand{\geq}{\geqslant}
\usepackage{polytable}

%mathindent has to be defined
\@ifundefined{mathindent}%
  {\newdimen\mathindent\mathindent\leftmargini}%
  {}%

\def\resethooks{%
  \global\let\SaveRestoreHook\empty
  \global\let\ColumnHook\empty}
\newcommand*{\savecolumns}[1][default]%
  {\g@addto@macro\SaveRestoreHook{\savecolumns[#1]}}
\newcommand*{\restorecolumns}[1][default]%
  {\g@addto@macro\SaveRestoreHook{\restorecolumns[#1]}}
\newcommand*{\aligncolumn}[2]%
  {\g@addto@macro\ColumnHook{\column{#1}{#2}}}

\resethooks

\newcommand{\onelinecommentchars}{\quad-{}- }
\newcommand{\commentbeginchars}{\enskip\{-}
\newcommand{\commentendchars}{-\}\enskip}

\newcommand{\visiblecomments}{%
  \let\onelinecomment=\onelinecommentchars
  \let\commentbegin=\commentbeginchars
  \let\commentend=\commentendchars}

\newcommand{\invisiblecomments}{%
  \let\onelinecomment=\empty
  \let\commentbegin=\empty
  \let\commentend=\empty}

\visiblecomments

\newlength{\blanklineskip}
\setlength{\blanklineskip}{0.66084ex}

\newcommand{\hsindent}[1]{\quad}% default is fixed indentation
\let\hspre\empty
\let\hspost\empty
\newcommand{\NB}{\textbf{NB}}
\newcommand{\Todo}[1]{$\langle$\textbf{To do:}~#1$\rangle$}

\EndFmtInput
\makeatother
%
%
%
%
%
%
% This package provides two environments suitable to take the place
% of hscode, called "plainhscode" and "arrayhscode". 
%
% The plain environment surrounds each code block by vertical space,
% and it uses \abovedisplayskip and \belowdisplayskip to get spacing
% similar to formulas. Note that if these dimensions are changed,
% the spacing around displayed math formulas changes as well.
% All code is indented using \leftskip.
%
% Changed 19.08.2004 to reflect changes in colorcode. Should work with
% CodeGroup.sty.
%
\ReadOnlyOnce{polycode.fmt}%
\makeatletter

\newcommand{\hsnewpar}[1]%
  {{\parskip=0pt\parindent=0pt\par\vskip #1\noindent}}

% can be used, for instance, to redefine the code size, by setting the
% command to \small or something alike
\newcommand{\hscodestyle}{}

% The command \sethscode can be used to switch the code formatting
% behaviour by mapping the hscode environment in the subst directive
% to a new LaTeX environment.

\newcommand{\sethscode}[1]%
  {\expandafter\let\expandafter\hscode\csname #1\endcsname
   \expandafter\let\expandafter\endhscode\csname end#1\endcsname}

% "compatibility" mode restores the non-polycode.fmt layout.

\newenvironment{compathscode}%
  {\par\noindent
   \advance\leftskip\mathindent
   \hscodestyle
   \let\\=\@normalcr
   \let\hspre\(\let\hspost\)%
   \pboxed}%
  {\endpboxed\)%
   \par\noindent
   \ignorespacesafterend}

\newcommand{\compaths}{\sethscode{compathscode}}

% "plain" mode is the proposed default.
% It should now work with \centering.
% This required some changes. The old version
% is still available for reference as oldplainhscode.

\newenvironment{plainhscode}%
  {\hsnewpar\abovedisplayskip
   \advance\leftskip\mathindent
   \hscodestyle
   \let\hspre\(\let\hspost\)%
   \pboxed}%
  {\endpboxed%
   \hsnewpar\belowdisplayskip
   \ignorespacesafterend}

\newenvironment{oldplainhscode}%
  {\hsnewpar\abovedisplayskip
   \advance\leftskip\mathindent
   \hscodestyle
   \let\\=\@normalcr
   \(\pboxed}%
  {\endpboxed\)%
   \hsnewpar\belowdisplayskip
   \ignorespacesafterend}

% Here, we make plainhscode the default environment.

\newcommand{\plainhs}{\sethscode{plainhscode}}
\newcommand{\oldplainhs}{\sethscode{oldplainhscode}}
\plainhs

% The arrayhscode is like plain, but makes use of polytable's
% parray environment which disallows page breaks in code blocks.

\newenvironment{arrayhscode}%
  {\hsnewpar\abovedisplayskip
   \advance\leftskip\mathindent
   \hscodestyle
   \let\\=\@normalcr
   \(\parray}%
  {\endparray\)%
   \hsnewpar\belowdisplayskip
   \ignorespacesafterend}

\newcommand{\arrayhs}{\sethscode{arrayhscode}}

% The mathhscode environment also makes use of polytable's parray 
% environment. It is supposed to be used only inside math mode 
% (I used it to typeset the type rules in my thesis).

\newenvironment{mathhscode}%
  {\parray}{\endparray}

\newcommand{\mathhs}{\sethscode{mathhscode}}

% texths is similar to mathhs, but works in text mode.

\newenvironment{texthscode}%
  {\(\parray}{\endparray\)}

\newcommand{\texths}{\sethscode{texthscode}}

% The framed environment places code in a framed box.

\def\codeframewidth{\arrayrulewidth}
\RequirePackage{calc}

\newenvironment{framedhscode}%
  {\parskip=\abovedisplayskip\par\noindent
   \hscodestyle
   \arrayrulewidth=\codeframewidth
   \tabular{@{}|p{\linewidth-2\arraycolsep-2\arrayrulewidth-2pt}|@{}}%
   \hline\framedhslinecorrect\\{-1.5ex}%
   \let\endoflinesave=\\
   \let\\=\@normalcr
   \(\pboxed}%
  {\endpboxed\)%
   \framedhslinecorrect\endoflinesave{.5ex}\hline
   \endtabular
   \parskip=\belowdisplayskip\par\noindent
   \ignorespacesafterend}

\newcommand{\framedhslinecorrect}[2]%
  {#1[#2]}

\newcommand{\framedhs}{\sethscode{framedhscode}}

% The inlinehscode environment is an experimental environment
% that can be used to typeset displayed code inline.

\newenvironment{inlinehscode}%
  {\(\def\column##1##2{}%
   \let\>\undefined\let\<\undefined\let\\\undefined
   \newcommand\>[1][]{}\newcommand\<[1][]{}\newcommand\\[1][]{}%
   \def\fromto##1##2##3{##3}%
   \def\nextline{}}{\) }%

\newcommand{\inlinehs}{\sethscode{inlinehscode}}

% The joincode environment is a separate environment that
% can be used to surround and thereby connect multiple code
% blocks.

\newenvironment{joincode}%
  {\let\orighscode=\hscode
   \let\origendhscode=\endhscode
   \def\endhscode{\def\hscode{\endgroup\def\@currenvir{hscode}\\}\begingroup}
   %\let\SaveRestoreHook=\empty
   %\let\ColumnHook=\empty
   %\let\resethooks=\empty
   \orighscode\def\hscode{\endgroup\def\@currenvir{hscode}}}%
  {\origendhscode
   \global\let\hscode=\orighscode
   \global\let\endhscode=\origendhscode}%

\makeatother
\EndFmtInput
%
%% ----------------------
%% Stuff for TimCore


%% 'g' inferences




%% --------- Effects ------------


%% ----------------------




%% ----------- Applications --------------
%% %format applyx a b = a "@" b
%% %format applyx a b = a "\," b

%% Function application with no arguments

%% Function application with 1 argument

%% Function application with 1 argument, without parenthesis

%% Function application with many arguments

%% ----------- End of Applications --------------


%% Sequences e1; ...; en; v

%% xs etc stole syntax used in examples.

%% Old version: | for BNF, [] for choice
%%%format choice = "[ \! ]"
%%%format `choice` = "\mathop{[ \! ]}"
%% 1mm is about 0.3em with the current font

%% New version: tall | for BNF, fat | for choice
%%format vbar = "\mathop{\bigm|}"
%%format choice = "\boldsymbol{|}"
%%format `choice` = "\mathop{\boldsymbol{|}}"




%% Arrays and tuples
%% %format arrKW = "\textbf{arr}"
%% %format arr (x) = arrKW "\{" x "\}"



%% %format defKW = "\textbf{def}"

% only used in comments, but still prettier this way:
%%%%format map = "\textbf{map}"

%% Effects checking, like succeeds{e}


%% Unification
%%%format == = "\!=\!"


% format := = "\mapsto"
%% %format squiggt = "\!\leadsto\!"
%  format effs x = "\!\!<\!\!" x "\!\!>\!\!" dots


































%% ----------- Metatheory --------------
%%%%%%%%%%%%%%%%%%%%%%%%%%%%%%%%%%%%

\newcommand\mydots{\makebox[0.55em][c]{\text{$\cdot\hspace{-0.3em}\cdot\hspace{-0.3em}\cdot$}}}


%% Complicated subscript tweak so that a subscript is the same height whether or not a superscript is present.
%% See Knuth, "The TeXBook", page 179, second "double-dangerous-turn" paragraph.
%% The implementation here affects only subscripts in the \mathcal font.
\newif\ifalignsubscripts

\global\alignsubscriptstrue

\newcommand{\setsubscriptalignparameters}{
\ifalignsubscripts
{\setbox0=\hbox{$\mathcal{A
\global\fontdimen16\textfont2=\fontdimen17\textfont2
}$}}
\fi
}

\setsubscriptalignparameters

\newcommand{\alignsubscript}[1]{\ifalignsubscripts\smash{#1}\else#1\fi}
%\newcommand{\alignsubscript}[1]{\ifalignsubscripts\vbox to 0pt{\vss$#1$}\else#1\fi}


% Typesetting rewrite rules
\newcommand{\rulename}[1]{\textsc{\small #1}}
\newcommand{\vb}{\; \ensuremath{\mathop{\mid}} \;}    % Vertical bar
\newcommand{\hole}{\square}
\newcommand{\freevars}[1]{\mathsf{fvs}(#1)}
\newcommand{\freevarsoutsidelambdas}[1]{\mathsf{fvsol}(#1)}
\newcommand{\boundvars}[1]{\mathsf{bvs}(#1)}
\newcommand{\valvars}[1]{\mathsf{vvs}(#1)}
\newcommand{\disjoint}{\;\#\;}
\newcommand{\hnf}{\mathit{hnf}}  % head-normal form
\newcommand{\op}{\mathit{op}}
\newcommand{\subst}[3]{#1\{#3/#2\}}

% Name of the language
\newcommand{\versecalc}{$\mathcal{VC}$}
\newcommand{\calcname}{\versecalc\xspace}
%%\newcommand{\mypara}[1]{\smallskip\noindent\textbf{#1}\xspace}  ACMART gives warning about use of \\smallskip for some reason
\newcommand{\mypara}[1]{\vskip\smallskipamount\noindent\textbf{#1}\xspace}

% \newcommand{\calcname}{Fresh Calculus\xspace}
% \newcommand{\lcalcname}{Lenient Calculus\xspace}
% \newcommand{\calcshort}{FC\xspace}
\newcommand{\vcname}{Verse Calculus\xspace}
\newcommand{\vcshort}{VC\xspace}

% Right arrow for rewrites
\newcommand{\movesto}{\longrightarrow}
\newcommand{\movestoX}{\rightsquigarrow}
\newcommand{\movesfromX}{\leftsquigarrow}
\newcommand{\movesfrom}{\longleftarrow}
\newcommand{\convertible}{\longleftrightarrow^{\ast}}

% Arrows with explanation, including the common cases of one, two or a repeated rule
\newcommand{\movestoC}[1]{\movesto\!\!\{#1\}}
\newcommand{\movestoR}[1]{\movestoC{\rulename{#1}}}
\newcommand{\movestoRR}[2]{\movestoC{\rulename{#1}, \rulename{#2}}}
\newcommand{\movestoRn}[2]{\movestoC{\rulename{#1} \times #2}}


\newcommand{\rightarrowdbl}{\rightarrow\mathrel{\mkern-14mu}\rightarrow}
\newcommand{\xrightarrowdbl}[2][]{%
  \xrightarrow[#1]{#2}\mathrel{\mkern-14mu}\rightarrow
}
\newcommand{\leftarrowdbl}{\leftarrow\mathrel{\mkern-14mu}\leftarrow}
\newcommand{\xleftarrowdbl}[2][]{%
  \leftarrow\mathrel{\mkern-14mu}\xleftarrow[#1]{#2}%
}

\newcommand{\xleftrightarrowdbl}[2][]{%
  \leftarrow\mathrel{\mkern-14mu}\leftarrow\mathrel{\mkern-14mu}\xrightarrow[#1]{#2}\mathrel{\mkern-14mu}\rightarrow}


\newcommand{\rr}{\ensuremath{\mathcal{R}}\xspace}
\newcommand{\rs}{\ensuremath{\mathcal{S}}\xspace}
\newcommand{\rt}{\ensuremath{\mathcal{T}}\xspace}
\newcommand{\rrhat}{\trel{\rr}}
\newcommand{\rshat}{\trel{\rs}}
\newcommand{\rthat}{\trel{\rt}}
\newcommand{\goone}[1]{\xrightarrow{\epsilon}_{\alignsubscript{#1}}}
\newcommand{\go}[1]{\xrightarrow{}_{\alignsubscript{#1}}}
\newcommand{\gos}[1]{\xrightarrowdbl{}_{\alignsubscript{#1}}}
\newcommand{\lgos}[1]{\xleftarrowdbl{}_{\alignsubscript{#1}}}
\newcommand{\gok}[2]{\xrightarrowdbl{#2}_{\alignsubscript{#1}}}

\newcommand{\trel}[1]{\ensuremath{\widehat{#1}}}
\newcommand{\tstep}[3]{\ensuremath{#2 \go{\trel{#1}} #3}}

\newcommand{\rstepo}[3]{\ensuremath{#2 \goone{#1} #3}}
\newcommand{\rstep}[3]{\ensuremath{#2 \go{#1} #3}}
\newcommand{\rsteps}[3]{\ensuremath{#2 \gos{#1} #3}}
\newcommand{\lsteps}[3]{\ensuremath{#2 \lgos{#1} #3}}
\newcommand{\rstepk}[4]{\ensuremath{#3 \gok{#1}{#2} #4}}
\newcommand{\ungoone}[1]{\xleftarrow{\epsilon}_{\alignsubscript{#1}}}
\newcommand{\ungo}[1]{\xleftarrow{}_{\alignsubscript{#1}}}
\newcommand{\ungos}[1]{\xleftarrowdbl{}_{\alignsubscript{#1}}}
\newcommand{\ungok}[2]{\xleftarrowdbl{#2}_{\alignsubscript{#1}}}
\newcommand{\unrstepo}[3]{\ensuremath{#2 \ungoone{#1} #3}}
\newcommand{\unrstep}[3]{\ensuremath{#2 \ungo{#1} #3}}
\newcommand{\unrsteps}[3]{\ensuremath{#2 \ungos{#1} #3}}
\newcommand{\unrstepk}[4]{\ensuremath{#3 \ungok{#1}{#2} #4}}
\newcommand{\univ}{\ensuremath{\mathcal{E}}}
\newcommand{\join}[3]{{#2} \downarrow_{\alignsubscript{#1}} {#3}}
\newcommand{\dia}{\ensuremath{\diamond}}

\newcommand{\rnk}{\ensuremath{\rho}}
\newcommand{\rord}{\ensuremath{\mathcal{O}}}
\newcommand{\runi}{\ensuremath{\mathcal{U}}}
\newcommand{\runihat}{\ensuremath{\trel{\mathcal{U}}}}
\newcommand{\rstr}{\ensuremath{\mathcal{S}}}
\newcommand{\rapp}{\ensuremath{\mathcal{A}}}
\newcommand{\rapphat}{\ensuremath{\trel{\mathcal{A}}}}
\newcommand{\rcon}{\ensuremath{\mathcal{N}}}
\newcommand{\rconhat}{\ensuremath{\trel{\mathcal{N}}}}
\newcommand{\rconm}{\ensuremath{\rcon'}}
\newcommand{\rss}{\ensuremath{\mathcal{SS}}}

\newcommand{\rgc}{\ensuremath{\mathcal{G}}}
\newcommand{\rgchat}{\ensuremath{\trel{\mathcal{G}}}}
\newcommand{\rfail}{\ensuremath{\mathcal{F}}}
\newcommand{\rspc}{\ensuremath{\mathcal{C}}}
\newcommand{\rspchat}{\ensuremath{\trel{\mathcal{C}}}}
\newcommand{\norecursion}{no recursion\xspace}
\newcommand{\noHYPHENrecursion}{no-recursion\xspace}
\newcommand{\rdx}[1]{\ensuremath{\Delta_{\alignsubscript{#1}}}}
\newcommand{\txr}[2]{\ensuremath{\rstep{#1}{\rdx{#2}}{\rdx{#2}'}}}
\newcommand{\tx}[1]{\ensuremath{\rstep{#1}{\rdx{#1}}{\rdx{#1}'}}}
\newcommand{\txi}[1]{\ensuremath{\rstep{}{\rdx{#1}}{\rdx{#1}'}}}


\newcommand{\wf}[1]{\ensuremath{\vdash #1}}
\newcommand{\wff}[5]{\ensuremath{#1 \vdash #2 \rightsquigarrow (#3 \mid #4 \mid #5)}}
\newcommand{\wffv}[4]{\ensuremath{\vdash_\ensuremath{\Varid{v}} #1 \rightsquigarrow (#2 \mid #3 \mid #4)}}
\newcommand{\wffmany}[4]{\ensuremath{\vdash_{\!*} #1 \rightsquigarrow (#2 \mid #3 \mid #4)}}
% \newcommand{\lvc}[1]{\colorbox{lightgray}{#1}}
% \newcommand{\lvc}[1]{\colorbox{lightgray}{#1}}
\newcommand{\lvc}[1]{#1}
\newcommand{\relDescription}[1]{\ensuremath{\textrm{\textbf{#1}}}}
\newcommand{\judgementHead}[2]{\ensuremath{\relDescription{#1}\hfill\fbox{#2}}}
\newcommand{\judgementHeadNameOnly}[1]{\ensuremath{\relDescription{#1}\hfill}}


%% Definitions for explaining skew confluence
\newcommand{\inccon}[2]{\mathbin{\smash{\mathop{\xleftrightarrowdbl{\hphantom{#1}}}\limits_{\smash{\hbox{\raise0.55ex\hbox{\smash{$\scriptstyle#2$}}}}}^{\vbox to 0pt{\hbox{\lower0.95ex\hbox{\smash{$\scriptstyle#1$}}}\vss}}}}}
\newcommand{\leqomega}{\mathrel{\leq_\omega}} 
\newcommand{\omegarr}{\omega_{\rr}}
\newcommand{\leqomegarr}{\mathrel{\leq_{\omegarr}}} 
\newcommand{\preceqomegarr}{\mathrel{\preceq_{\omega_{\rr}}}}
\newcommand{\downclose}{\mathord{\Downarrow_{\leqomega}\!}}


%% GS: Add @@{} to match other changes made below to get more horizontal space in Figure 3
%% \newcommand{\rulegroupname}[1]{\multicolumn{4}{l}{\textit{#1}}\\}
\newcommand{\rulegroupname}[1]{\multicolumn{4}{@{}l}{\textit{#1}}\\}
\newcommand{\largerhs}[1]{\multicolumn{2}{@{}l}{#1}}



%  -------------------------- Syntax -----------------------
\newcommand{\figureSyntax}{
\begin{figure}[t]
    $$
%%  \begin{array}{l}
  \begin{array}{l@{}}
      \textbf{Syntax} \\
      \begin{array}{lr@{\;}c@{\;}l@{}}
      \text{Integers} & k \\
      \text{Variables} & \multicolumn{3}{l}{x,y,z,f,g} \\
      \text{Programs} & p  & ::=  & \ensuremath{\textbf{one}\{\Varid{e}\}} \quad \text{where} \; \freevars{e} = \emptyset \\
      \text{Expressions}   & e     & ::= & \ensuremath{\Varid{v}} \vb \ensuremath{eq;\,\Varid{e}} \vb \ensuremath{\exists\Varid{x}.\,\Varid{e}} \vb \ensuremath{\textbf{fail}}
                                           \vb \ensuremath{\Varid{e}_{\mathrm{1}}\;\rule[-0.1em]{0.15em}{0.75em}\;\Varid{e}_{\mathrm{2}}} \vb \ensuremath{\Varid{v}_{\mathrm{1}}\,\Varid{v}_{\mathrm{2}}}
                             \vb \ensuremath{\textbf{one}\{\Varid{e}\}} \vb \ensuremath{\textbf{all}\{\Varid{e}\}} \\
                             & \ensuremath{eq}     & ::= & \ensuremath{\Varid{e}} \vb \ensuremath{\Varid{v}\hspace{-0.14em}=\hspace{-0.14em}\Varid{e}} \quad \quad \quad \quad
                                       \text{Note: ``\ensuremath{eq}'' is pronounced ``expression or equation''} \\
  %    \text{Expressions} & e  & ::=  & |v| \vb |apply1N v_1 v_2|
  %       \vb |e_1 == e_2| \vb |e_1 ; e_2| \vb  |def x e|
  %    \vb  |e_1 choice e_2| \vb |fail| \vb |one e| \vb |all e| \\
  % SPJ Omitting for now
  \iffullverse
      && \vb & \ensuremath{\textbf{succeeds}\;\Varid{e}} \vb \ensuremath{\textbf{decides}\;\Varid{e}} \\
  \fi
      \text{Values} & \ensuremath{\Varid{v}} & ::= & x  \vb \hnf \\
      \text{Head values} & \hnf & ::= & k \vb op \vb \ensuremath{\langle\Varid{v}_{\mathrm{1}},\mydots,\Varid{v}_{\Varid{n}}\rangle} \vb \ensuremath{\lambda\Varid{x}.\,\Varid{e}} \\
      \text{Primops} & op  & ::= & \ensuremath{\textbf{gt}} \vb \ensuremath{\textbf{add}}
  \iffullverse
      \vb \ensuremath{\textbf{isInt}}
  \fi
      \end{array} \\
      \\[-0.9em]
      \text{\emph{Concrete syntax}:} \;\;
      \begin{array}[t]{l}
            \text{``\ensuremath{\rule[-0.1em]{0.15em}{0.75em}}'' and ``;'' are right-associative} \\
            \text{``\ensuremath{\rule[-0.1em]{0.15em}{0.75em}}'' and ``\ensuremath{\hspace{-0.14em}=\hspace{-0.14em}}'' bind more tightly than ``;''} \\
            \text{``$\lambda$'', ``\ensuremath{\exists{}}'' scope as far to the right as possible} \\
            \hspace{1.2em} \text{{\eg}, \ensuremath{(\lambda\Varid{y}.\,\exists\Varid{x}.\,\Varid{x}\hspace{-0.14em}=\hspace{-0.14em}\mathrm{1};\,\Varid{x}\mathbin{+}\Varid{y})} means \ensuremath{(\lambda\Varid{y}.\,(\exists\Varid{x}.\,((\Varid{x}\hspace{-0.14em}=\hspace{-0.14em}\mathrm{1});\,(\Varid{x}\mathbin{+}\Varid{y}))))}} \\
      \end{array} \\
      \text{Parentheses may be used freely to aid readability and override default precedence.}\\
  \\[-0.6em]
      \textbf{Desugaring} \\
%%    \begin{array}{rcl@@{\hspace{4em}}l}
    \begin{array}{rcl@{\hspace{3.2em}}l}
%%      |emptytup| & \text{means} & |arr ()|\\
%%      |tup (v_1, dots, v_n)| & \text{means} & |arr (v_1,dots, v_n)| \quad n \geq 2 \\
\iffalse
% Do not use := when defining sugar
      \ensuremath{\Varid{v}_{\mathrm{1}}\mathbin{+}\Varid{v}_{\mathrm{2}}} & \text{means} & \ensuremath{\textbf{add}\langle\Varid{v}_{\mathrm{1}},\Varid{v}_{\mathrm{2}}\rangle} & \\
      \ensuremath{\Varid{v}_{\mathrm{1}}\mathbin{>}\Varid{v}_{\mathrm{2}}} & \text{means} & \ensuremath{\textbf{gt}\langle\Varid{v}_{\mathrm{1}},\Varid{v}_{\mathrm{2}}\rangle} & \\
      \ensuremath{\exists\Varid{x}_{\mathrm{1}}\hspace{0.1em}\Varid{x}_{\mathrm{2}}\;\mydots\;\Varid{x}_{\Varid{n}}.\,\Varid{e}} & \text{means} & \ensuremath{\exists\Varid{x}_{\mathrm{1}}.\,\exists\Varid{x}_{\mathrm{2}}.\,\mydots \exists\Varid{x}_{\Varid{n}}.\,\Varid{e}} & \\
      \ensuremath{\Varid{e}_{\mathrm{1}}\;\Varid{e}_{\mathrm{2}}} & \text{means}^* & \ensuremath{\exists\Varid{f}\hspace{0.1em}\Varid{x}.\,\Varid{f}\hspace{-0.14em}=\hspace{-0.14em}\Varid{e}_{\mathrm{1}};\,\Varid{x}\hspace{-0.14em}=\hspace{-0.14em}\Varid{e}_{\mathrm{2}};\,\Varid{f}\,\Varid{x}}  & \ensuremath{\Varid{f},\Varid{x}}\ \text{fresh} \\
      \ensuremath{\langle{}\Varid{e}_{\mathrm{1}},\mydots,\Varid{e}_{\Varid{n}}\rangle} & \text{means}^* &
          \ensuremath{\exists\Varid{x}_{\mathrm{1}}\hspace{0.1em}\Varid{x}_{\mathrm{2}}\;\mydots\;\Varid{x}_{\Varid{n}}.\,\Varid{x}_{\mathrm{1}}\hspace{-0.14em}=\hspace{-0.14em}\Varid{e}_{\mathrm{1}};\,\mydots;\,\Varid{x}_{\Varid{n}}\hspace{-0.14em}=\hspace{-0.14em}\Varid{e}_{\Varid{n}};\,\langle{}\Varid{x}_{\mathrm{1}},\mydots,\Varid{x}_{\Varid{n}}\rangle} & \text{\ensuremath{\Varid{x}_{\Varid{i}}} fresh} \\
      \ensuremath{\Varid{e}_{\mathrm{1}}\hspace{-0.14em}=\hspace{-0.14em}\Varid{e}_{\mathrm{2}}}      & \text{means}^* & \ensuremath{\exists\Varid{x}.\,\Varid{x}\hspace{-0.14em}=\hspace{-0.14em}\Varid{e}_{\mathrm{1}};\,\Varid{x}\hspace{-0.14em}=\hspace{-0.14em}\Varid{e}_{\mathrm{2}};\,\Varid{x}}       & \ensuremath{\Varid{x}}\ \text{fresh} \\
%      |e_1 ; e_2|      & \text{means} & |def x (e_1; x == e_2 ; x)|       & |x|\ \text{fresh} \\

%       |for^(e_1) do e_2| & \text{means} & |def v (v == all (e_1; thunk e_2); applyN map (lam z (apply0 z), v))|
      \ensuremath{\Varid{x}\hspace{-0.14em}:=\hspace{-0.14em}\Varid{e}_{\mathrm{1}};\,\Varid{e}_{\mathrm{2}}} & \text{means} & \ensuremath{\exists\Varid{x}.\,\Varid{x}\hspace{-0.14em}=\hspace{-0.14em}\Varid{e}_{\mathrm{1}};\,\Varid{e}_{\mathrm{2}}} \\
      \ensuremath{\lambda\langle\rangle.\,\hspace{0.1em}\Varid{e}} & \text{means} & \ensuremath{\lambda\Varid{p}.\,\Varid{p}\hspace{-0.14em}=\hspace{-0.14em}\langle\rangle;\,\Varid{e}} & \text{\ensuremath{\Varid{p}} fresh} \\
      \ensuremath{\lambda\langle{}\Varid{x},\Varid{y}\rangle.\,\hspace{0.1em}\Varid{e}} & \text{means} & \ensuremath{\lambda\Varid{p}.\,\exists\Varid{x}\hspace{0.1em}\Varid{y}.\,\Varid{p}\hspace{-0.14em}=\hspace{-0.14em}\langle{}\Varid{x},\Varid{y}\rangle;\,\Varid{e}} & \text{\ensuremath{\Varid{p}} fresh} \\
      \ensuremath{\textbf{if}\;\Varid{e}_{\mathrm{1}}\;\textbf{then}\;\Varid{e}_{\mathrm{2}}\;\textbf{else}\;\Varid{e}_{\mathrm{3}}} & \text{means}
      & \ensuremath{\exists\Varid{y}.\,\Varid{y}\hspace{-0.14em}=\hspace{-0.14em}\textbf{one}\{(\Varid{e}_{\mathrm{1}};\,\lambda \langle\rangle.\,\Varid{e}_{\mathrm{2}})\;\rule[-0.1em]{0.15em}{0.75em}\;(\lambda \langle\rangle.\,\Varid{e}_{\mathrm{3}})\};\,\Varid{y}\langle\rangle} & \text{\ensuremath{\Varid{y}} fresh} \\
\else
% Use := when defining sugar
      \ensuremath{\Varid{e}_{\mathrm{1}}\mathbin{+}\Varid{e}_{\mathrm{2}}} & \text{means} & \ensuremath{\textbf{add}\langle\Varid{e}_{\mathrm{1}},\Varid{e}_{\mathrm{2}}\rangle} & \\
      \ensuremath{\Varid{e}_{\mathrm{1}}\mathbin{>}\Varid{e}_{\mathrm{2}}} & \text{means} & \ensuremath{\textbf{gt}\langle\Varid{e}_{\mathrm{1}},\Varid{e}_{\mathrm{2}}\rangle} & \\
      \ensuremath{\exists\Varid{x}_{\mathrm{1}}\hspace{0.1em}\Varid{x}_{\mathrm{2}}\;\mydots\;\Varid{x}_{\Varid{n}}.\,\Varid{e}} & \text{means} & \ensuremath{\exists\Varid{x}_{\mathrm{1}}.\,\exists\Varid{x}_{\mathrm{2}}.\,\mydots \exists\Varid{x}_{\Varid{n}}.\,\Varid{e}} & \\
      \ensuremath{\Varid{x}\hspace{-0.14em}:=\hspace{-0.14em}\Varid{e}_{\mathrm{1}};\,\Varid{e}_{\mathrm{2}}} & \text{means} & \ensuremath{\exists\Varid{x}.\,\Varid{x}\hspace{-0.14em}=\hspace{-0.14em}\Varid{e}_{\mathrm{1}};\,\Varid{e}_{\mathrm{2}}} \\
      \ensuremath{\Varid{e}_{\mathrm{1}}\;\Varid{e}_{\mathrm{2}}} & \text{means}^* & \ensuremath{\Varid{f}\hspace{-0.14em}:=\hspace{-0.14em}\Varid{e}_{\mathrm{1}};\,\Varid{x}\hspace{-0.14em}:=\hspace{-0.14em}\Varid{e}_{\mathrm{2}};\,\Varid{f}\,\Varid{x}}  & \ensuremath{\Varid{f},\Varid{x}}\ \text{fresh} \\
      \ensuremath{\langle{}\Varid{e}_{\mathrm{1}},\mydots,\Varid{e}_{\Varid{n}}\rangle} & \text{means}^* &
          \ensuremath{\Varid{x}_{\mathrm{1}}\hspace{-0.14em}:=\hspace{-0.14em}\Varid{e}_{\mathrm{1}};\,\mydots;\,\Varid{x}_{\Varid{n}}\hspace{-0.14em}:=\hspace{-0.14em}\Varid{e}_{\Varid{n}};\,\langle{}\Varid{x}_{\mathrm{1}},\mydots,\Varid{x}_{\Varid{n}}\rangle} & \text{\ensuremath{\Varid{x}_{\Varid{i}}} fresh} \\
      \ensuremath{\Varid{e}_{\mathrm{1}}\hspace{-0.14em}=\hspace{-0.14em}\Varid{e}_{\mathrm{2}}}      & \text{means}^* & \ensuremath{\Varid{x}\hspace{-0.14em}:=\hspace{-0.14em}\Varid{e}_{\mathrm{1}};\,\Varid{x}\hspace{-0.14em}=\hspace{-0.14em}\Varid{e}_{\mathrm{2}};\,\Varid{x}}       & \ensuremath{\Varid{x}}\ \text{fresh} \\

%       |for^(e_1) do e_2| & \text{means} & |def v (v == all (e_1; thunk e_2); applyN map (lam z (apply0 z), v))|
%      |lamhdr emptytup^^ e| & \text{means} & |lam p (p == emptytup; e)| & \text{|p| fresh} \\
%      |lamhdr (tup(x,y))^^ e| & \text{means} & |lam p (def (x ^^ y) (p == tup (x,y); e))| & \text{|p| fresh} \\
      \ensuremath{\lambda\langle{}\Varid{x}_{\mathrm{1}},\mydots,\Varid{x}_{\Varid{n}}\rangle.\,\hspace{0.1em}\Varid{e}} & \text{means} & \ensuremath{\lambda\Varid{p}.\,\exists\Varid{x}_{\mathrm{1}}\;\mydots\;\Varid{x}_{\Varid{n}}.\,\Varid{p}\hspace{-0.14em}=\hspace{-0.14em}\langle{}\Varid{x}_{\mathrm{1}},\mydots,\Varid{x}_{\Varid{n}}\rangle;\,\Varid{e}} & \text{\ensuremath{\Varid{p}} fresh, $\ensuremath{\Varid{n}} \geq 0$} \\
%      |if e_1 then e_2 else e_3| & \text{means}
%         & |y :== one( (e_1; thunk e_2) choice (thunk e_3)); apply0 y| & \text{|y| fresh} \\
      \ensuremath{\textbf{if}\;(\exists\Varid{x}_{\mathrm{1}}\hspace{0.1em}\mydots\hspace{0.1em}\Varid{x}_{\Varid{n}}.\,\Varid{e}_{\mathrm{1}})\;\textbf{then}\;\Varid{e}_{\mathrm{2}}\;\textbf{else}\;\Varid{e}_{\mathrm{3}}} & \text{means}
         & \multicolumn{2}{@{}l}{\ensuremath{(\textbf{one}\{(\exists\Varid{x}_{\mathrm{1}}\hspace{0.1em}\mydots\hspace{0.1em}\Varid{x}_{\Varid{n}}.\,\Varid{e}_{\mathrm{1}};\,\lambda \langle\rangle.\,\Varid{e}_{\mathrm{2}})\;\rule[-0.1em]{0.15em}{0.75em}\;(\lambda \langle\rangle.\,\Varid{e}_{\mathrm{3}})\})\langle\rangle}} \\
\fi
    \end{array} \\
  \\[-0.6em]
  (*)\; \text{Apply this rule only if at least one of the \ensuremath{\Varid{e}_{\Varid{i}}} is not a value \ensuremath{\Varid{v}}} \\
  \text{$\freevars{\ensuremath{\Varid{e}}}$ means the free variables of \ensuremath{\Varid{e}}; in \versecalc{}, $\lambda$ and \ensuremath{\exists{}} are the only binders.}
  %\text{$S_1 \disjoint S_2$ means that set $S_1$ is disjoint from $S_2$.}
  \end{array}
  $$
  \caption{\textbf{\calcname: Syntax}} \label{fig:syntax}
  \Description{A table that exhibits a BNF grammar for the syntax of the Verse Calculus, some rules about resolving parsing ambiguity,
      and a list of rules for desugaring constructs that are notationally convenient but not part of the Verse Calculus proper}
\end{figure}
}

%  -------------------------- Contexts -----------------------
\newcommand{\figureContexts}{
\begin{figure}
    $$
\begin{array}{l}
      \begin{array}{lr@{\;}c@{\;}l}
%      \text{Substitution contexts} & S & ::= & \hole \vb |eq; S| \\
      \text{Execution contexts} & X & ::= & \hole \vb \ensuremath{\Varid{v}\hspace{-0.14em}=\hspace{-0.14em}X;\,\Varid{e}} \vb \ensuremath{X;\,\Varid{e}} \vb \ensuremath{eq;\,X} \\
%       \text{Equation contexts (hole is |eq|)} & U & ::= & |hole ; e| \vb |v == U ; e| \vb |U ; e| \vb |e; U| \\
      \text{Value contexts} & V & ::= & \hole \vb \ensuremath{\langle{}\Varid{v}_{\mathrm{1}},\mydots,\Conid{V},\mydots,\Varid{v}_{\Varid{n}}\rangle}\\
      \text{Scope contexts} & SX & ::= & \ensuremath{\textbf{one}\{\Conid{SC}\}} \vb \ensuremath{\textbf{all}\{\Conid{SC}\}} \\
                            & SC & ::= & \ensuremath{\hole} \vb \ensuremath{\Conid{SC}\;\rule[-0.1em]{0.15em}{0.75em}\;\Varid{e}} \vb \ensuremath{\Varid{e}\;\rule[-0.1em]{0.15em}{0.75em}\;\Conid{SC}} \\

  \iffullverse
      && \vb & \ensuremath{\textbf{succeeds}\;\hole} \vb \ensuremath{\textbf{decides}\;\hole} \\
  \fi
      \text{Choice contexts} & CX & ::= & \hole \vb \ensuremath{\Varid{v}\hspace{-0.14em}=\hspace{-0.14em}CX;\,\Varid{e}} \vb \ensuremath{CX;\,\Varid{e}} \vb \ensuremath{\Varid{ceq};\,CX} \vb \ensuremath{\exists\Varid{x}.\,CX} \\
      \text{Choice-free exprs} & ce  & ::=  & \ensuremath{\Varid{v}} \vb \ensuremath{\Varid{ceq};\,\Varid{ce}}
                                     \vb \ensuremath{\textbf{one}\{\Varid{e}\}} \vb \ensuremath{\textbf{all}\{\Varid{e}\}} \vb \ensuremath{\exists\Varid{x}.\,\Varid{ce}}
                                     \vb \ensuremath{\Varid{op}(\Varid{v})} \\
                               & ceq & ::= & \ensuremath{\Varid{ce}} \vb \ensuremath{\Varid{v}\hspace{-0.14em}=\hspace{-0.14em}\Varid{ce}} \\
    \end{array}\\
\\[-0.6em]
\text{\textit{Note}: The \ensuremath{\hole} in \ensuremath{X} can only be an expression, not an equation.}
\end{array}
\vspace{-0.5em}
     $$
    \caption{The syntax of contexts}
    \Description{A table that exhibits a BNF grammar that defines execution contexts, value contexts., scope contexts, choice contexts, and chouce-free expressions}
      \label{fig:contexts} \label{fig:choice-free-syntax}
\end{figure}
}

\newcommand{\alphamark}{use $\alpha$}

% ------------------ ICFP rewrite rules -------------------

\newcommand{\figureRewriteRules}{
  %% BE SURE TO KEEP THE BULLET ITEMS IN \cref{sec:confluence} IN FILE ../confluence/04-confluence.ltx SYNCHRONIZED WITH THE RULE NAMES IN THIS FIGURE!
  \begin{figure}
%% GS: Add @@{} to match other changes made below to get more horizontal space in Figure 3 
%%    $\begin{array}{l}

         %%          $\begin{array}{@@{}l@@{}}
    $\begin{array}{@{}l}

    % \rulegroupname{Primitive operations}

    \iffullverse
        \rulename{p-isInt} & \ensuremath{\textbf{isInt}(\mathit{hnf})} & \movesto &
        $\left\{\begin{array}{@{}ll}
          \ensuremath{()} & \text{if \ensuremath{\mathit{hnf}} is an integer} \\
          \ensuremath{\textbf{fail}} & \text{otherwise}
          \end{array}\right.$
    \fi

%% GS: change \quad to \hskip0.5em to get a little more horizontal space to work with in this figure; also add @@{}
%% \begin{array}{@@{\quad} l  r @@{\;}c@@{\;} l l}
%%\begin{array}{@@{\hskip0.5em} l @@{}r @@{\;}c@@{\;} l@@{\hskip4em}l}

%% \begin{array}{@@{\hskip0.5em} l @@{\hspace{-.5em}} r @@{\;}c@@{\;} l l}
\begin{array}{@{\hskip0.5em} l @{\hspace{-.5em}} r @{\;}c@{\;} l l@{}}
\rulegroupname{Application:}
\rulename{app-add}          & \ensuremath{\textbf{add}\langle\Varid{k}_{\mathrm{1}},\Varid{k}_{\mathrm{2}}\rangle} & \movesto & \ensuremath{\Varid{k}_{\mathrm{3}}} & \text{where $k_3 = k_1 + k_2$}\\
\rulename{app-gt}           & \ensuremath{\textbf{gt}\langle\Varid{k}_{\mathrm{1}},\Varid{k}_{\mathrm{2}}\rangle}   & \movesto & \ensuremath{\Varid{k}_{\mathrm{1}}} &  \text{if $k_1 > k_2$} \\
\rulename{app-gt-fail}      & \ensuremath{\textbf{gt}\langle\Varid{k}_{\mathrm{1}},\Varid{k}_{\mathrm{2}}\rangle}   & \movesto & \ensuremath{\textbf{fail}} &  \text{if $k_1 \leq k_2$} \\
\rulename{app-beta}^{\alpha} & \ensuremath{(\lambda\Varid{x}.\,\Varid{e})(\Varid{v})} & \movesto & \ensuremath{\exists\Varid{x}.\,\Varid{x}\hspace{-0.14em}=\hspace{-0.14em}\Varid{v};\,\Varid{e}}
                                                  & \text{if $x \not\in \freevars{\ensuremath{\Varid{v}}}$} \\
\rulename{app-tup}          & \ensuremath{\langle{}\Varid{v}_{\mathrm{0}},\mydots,\Varid{v}_{\Varid{n}}\rangle(\Varid{v})} & \movesto &
                              \multicolumn{2}{@{}l}{\ensuremath{\exists\Varid{x}.\,\Varid{x}\hspace{-0.14em}=\hspace{-0.14em}\Varid{v};(\Varid{x}\hspace{-0.14em}=\hspace{-0.14em}\mathrm{0};\Varid{v}_{\mathrm{0}})\;\rule[-0.1em]{0.15em}{0.75em}\;\mydots\;\rule[-0.1em]{0.15em}{0.75em}\;(\Varid{x}\hspace{-0.14em}=\hspace{-0.14em}\Varid{n};\Varid{v}_{\Varid{n}})}
%%                              \hspace{1em} \text{fresh $|x| \not\in \freevars{|v, v_0, xdots,v_n|}$}} \\
                              \hspace{0.7em}\text{fresh $\ensuremath{\Varid{x}} \not\in \freevars{\ensuremath{\Varid{v},\Varid{v}_{\mathrm{0}},\mydots,\Varid{v}_{\Varid{n}}}}$}} \hskip-0.7em \\
\rulename{app-tup-0}        & \ensuremath{\langle\rangle(\Varid{v})} & \movesto & \ensuremath{\textbf{fail}} \\
%%\end{array}
%%\\
\\[-0.6em]

%% GS: change \quad to \hskip0.5em to get a little more horizontal space to work with in this figure 
%% \begin{array}{@@{\quad} l @@{\hspace{-.5em}} r @@{\;}c@@{\;} l l}
%%\begin{array}{@@{\hskip0.5em} l @@{\hspace{-.5em}} r @@{\;}c@@{\;} l l}
    % The quad indents the lines
    % The \hspace{-..} removes redundant space between rule name column and LHS column
    % that arises because the longest rule name and the longer LHS are not on the same line
    % The \; reduces the spacing around the arrow
\rulegroupname{Unification:}
  \rulename{u-lit}   & \ensuremath{\Varid{k}_{\mathrm{1}}\hspace{-0.14em}=\hspace{-0.14em}\Varid{k}_{\mathrm{2}};\,\Varid{e}} & \movesto & \ensuremath{\Varid{e}} & \text{if $k_1=k_2$} \\
%%  \rulename{u-tup} & \hspace{-1.8em} |tup (v_1,xdots,v_n) == tup (v1',xdots,v_n'); e| & \movesto & |v_1==v1';xdots;v_n==v_n'; e| \\
  \rulename{u-tup} & \hspace{-2.0em} \ensuremath{\langle{}\Varid{v}_{\mathrm{1}},\mydots,\Varid{v}_{\Varid{n}}\rangle\hspace{-0.14em}=\hspace{-0.14em}\langle{}\Varid{v}_{1}',\mydots,\Varid{v}_{\Varid{n}}'\rangle;\,\Varid{e}} & \movesto & \ensuremath{\Varid{v}_{\mathrm{1}}\hspace{-0.14em}=\hspace{-0.14em}\Varid{v}_{1}';\,\mydots;\,\Varid{v}_{\Varid{n}}\hspace{-0.14em}=\hspace{-0.14em}\Varid{v}_{\Varid{n}}';\,\Varid{e}} \\
  \rulename{u-fail} & \ensuremath{\mathit{hnf}_{\!1}\hspace{-0.14em}=\hspace{-0.14em}\mathit{hnf}_{\!2};\,\Varid{e}} & \movesto & \ensuremath{\textbf{fail}} & \text{if \rulename{u-lit}, \rulename{u-tup} do not match} \\
  \rulename{u-occurs} & \ensuremath{\Varid{x}\hspace{-0.14em}=\hspace{-0.14em}\Conid{V} [\mskip1.5mu \Varid{x}\mskip1.5mu];\,\Varid{e}} & \movesto & \ensuremath{\textbf{fail}} & \text{if $\ensuremath{\Conid{V}} \ne \ensuremath{\hole}$} \\
  \rulename{subst} & \ensuremath{X [\mskip1.5mu \Varid{x}\hspace{-0.14em}=\hspace{-0.14em}\Varid{v};\,\Varid{e}\mskip1.5mu]} & \movesto & \ensuremath{(\subst{X}{\Varid{x}}{\Varid{v}}) [\mskip1.5mu \Varid{x}\hspace{-0.14em}=\hspace{-0.14em}\Varid{v};\,\subst{\Varid{e}}{\Varid{x}}{\Varid{v}}\mskip1.5mu]} \hspace*{-0.3em}
  & \text{if $\ensuremath{\Varid{v}} \not= \ensuremath{\Conid{V} [\mskip1.5mu \Varid{x}\mskip1.5mu]}$} \\
  \rulename{hnf-swap} & \ensuremath{\mathit{hnf}\hspace{-0.14em}=\hspace{-0.14em}\Varid{x};\,\Varid{e}}    & \movesto & \ensuremath{\Varid{x}\hspace{-0.14em}=\hspace{-0.14em}\mathit{hnf};\,\Varid{e}} \\
  \rulename{var-swap} & \ensuremath{\Varid{y}\hspace{-0.14em}=\hspace{-0.14em}\Varid{x};\,\Varid{e}}    & \movesto & \ensuremath{\Varid{x}\hspace{-0.14em}=\hspace{-0.14em}\Varid{y};\,\Varid{e}} & \text{if \ensuremath{\Varid{x}} $\prec$ \ensuremath{\Varid{y}} } \\
%%  \rulename{var-swap-subst} & |X^[x==y; e]| & \movesto & |subst X y x^[y==x; subst e y x]| & \text{if |y| $\prec$ |x| } \\
%%  \rulename{val-swap-ord} & |eq_1; eq_2; e| & \movesto & |eq_2; eq_1; e| & \text{if |eq_2| $\prec$ |eq_1| } \\
  \rulename{seq-swap} & \ensuremath{eq;\,\Varid{x}\hspace{-0.14em}=\hspace{-0.14em}\Varid{v};\,\Varid{e}} & \movesto & \ensuremath{\Varid{x}\hspace{-0.14em}=\hspace{-0.14em}\Varid{v};\,eq;\,\Varid{e}} &
                      \text{unless (\ensuremath{eq} is \ensuremath{\Varid{y}\hspace{-0.14em}=\hspace{-0.14em}\Varid{v'}} and \ensuremath{\Varid{y}} $\preceq$ \ensuremath{\Varid{x}})} \\

    \\[-0.6em]

\rulegroupname{Elimination:}
\rulename{val-elim} & \ensuremath{\Varid{v};\,\Varid{e}} & \movesto & \ensuremath{\Varid{e}} \\
\rulename{exi-elim} & \ensuremath{\exists\Varid{x}.\,\Varid{e}} & \movesto & \ensuremath{\Varid{e}} & \text{if $\ensuremath{\Varid{x}} \notin \freevars{\ensuremath{\Varid{e}}}$}\\
\rulename{eqn-elim} & \ensuremath{\exists\Varid{x}.\,X [\mskip1.5mu \Varid{x}\hspace{-0.14em}=\hspace{-0.14em}\Varid{v};\,\Varid{e}\mskip1.5mu]} & \movesto & \ensuremath{X [\mskip1.5mu \Varid{e}\mskip1.5mu]} & \text{if $\ensuremath{\Varid{x}} \notin \freevars{\ensuremath{X [\mskip1.5mu \Varid{e}\mskip1.5mu]}}$ and $\ensuremath{\Varid{v}} \ne \ensuremath{\Conid{V} [\mskip1.5mu \Varid{x}\mskip1.5mu]}$} \\
\rulename{fail-elim} & \ensuremath{X [\mskip1.5mu \textbf{fail}\mskip1.5mu]} & \movesto & \ensuremath{\textbf{fail}} & \\

    \\[-0.6em]

\rulegroupname{Normalization:}
\rulename{exi-float}^{\alpha} & \ensuremath{X [\mskip1.5mu \exists\Varid{x}.\,\Varid{e}\mskip1.5mu]} & \movesto & \ensuremath{\exists\Varid{x}.\,X [\mskip1.5mu \Varid{e}\mskip1.5mu]} & \text{if $x \notin \freevars{\ensuremath{X}}$} \\
\rulename{seq-assoc} & \ensuremath{(eq;\,\Varid{e}_{\mathrm{1}});\,\Varid{e}_{\mathrm{2}}} & \movesto & \ensuremath{eq;\,(\Varid{e}_{\mathrm{1}};\,\Varid{e}_{\mathrm{2}})} \\
\rulename{eqn-float} & \ensuremath{\Varid{v}\hspace{-0.14em}=\hspace{-0.14em}(eq;\,\Varid{e}_{\mathrm{1}});\,\Varid{e}_{\mathrm{2}}} & \movesto & \ensuremath{eq;\,(\Varid{v}\hspace{-0.14em}=\hspace{-0.14em}\Varid{e}_{\mathrm{1}};\,\Varid{e}_{\mathrm{2}})} \\
\rulename{exi-swap} & \ensuremath{\exists\Varid{x}.\,\exists\Varid{y}.\,\Varid{e}} & \movesto & \ensuremath{\exists\Varid{y}.\,\exists\Varid{x}.\,\Varid{e}} \\
%\rulename{seq-swap} & |eq_1; eq_2; e| & \movesto & |eq_2; eq_1; e| &
%                      \hspace{-3em}\text{if ($eq_1 \in ceq$ or $eq_2 \in ceq$) and $eq_2 \prec eq_1$} \\
%\multicolumn{5}{l}{\text{\lennart{Plan: remove rule above, use rule below, move to unification group}}}\\

    \\[-0.6em]

\rulegroupname{Choice:}
\rulename{one-fail}   & \ensuremath{\textbf{one}\{\textbf{fail}\}} & \movesto & \ensuremath{\textbf{fail}} \\
\rulename{one-value}  & \ensuremath{\textbf{one}\{\Varid{v}\}} & \movesto & \ensuremath{\Varid{v}} \\
\rulename{one-choice} & \ensuremath{\textbf{one}\{\Varid{v}\hspace{0.2em}\mathop{\rule[-0.1em]{0.15em}{0.75em}}\hspace{0.2em}\Varid{e}\}} & \movesto & \ensuremath{\Varid{v}} \\
\rulename{all-fail}   & \ensuremath{\textbf{all}\{\textbf{fail}\}} & \movesto & \ensuremath{\langle{}\rangle} \\
\rulename{all-value}  & \ensuremath{\textbf{all}\{\Varid{v}\}} & \movesto & \ensuremath{\langle{}\Varid{v}\rangle} \\
\rulename{all-choice} & \ensuremath{\textbf{all}\{\Varid{v}_{\mathrm{1}}\;\rule[-0.1em]{0.15em}{0.75em}\;\mydots\;\rule[-0.1em]{0.15em}{0.75em}\;\Varid{v}_{\Varid{n}}\}} & \movesto & \ensuremath{\langle{}\Varid{v}_{\mathrm{1}},\mydots,\Varid{v}_{\Varid{n}}\rangle} \\
\rulename{choose-r} & \ensuremath{\textbf{fail}\;\rule[-0.1em]{0.15em}{0.75em}\;\Varid{e}} & \movesto & \ensuremath{\Varid{e}} \\
\rulename{choose-l} & \ensuremath{\Varid{e}\;\rule[-0.1em]{0.15em}{0.75em}\;\textbf{fail}} & \movesto & \ensuremath{\Varid{e}} \\
\rulename{choose-assoc} & \ensuremath{(\Varid{e}_{\mathrm{1}}\;\rule[-0.1em]{0.15em}{0.75em}\;\Varid{e}_{\mathrm{2}})\;\rule[-0.1em]{0.15em}{0.75em}\;\Varid{e}_{\mathrm{3}}} & \movesto & \ensuremath{\Varid{e}_{\mathrm{1}}\;\rule[-0.1em]{0.15em}{0.75em}\;(\Varid{e}_{\mathrm{2}}\;\rule[-0.1em]{0.15em}{0.75em}\;\Varid{e}_{\mathrm{3}})} \\
\rulename{choose} & \ensuremath{\Conid{SX} [\mskip1.5mu CX [\mskip1.5mu \Varid{e}_{\mathrm{1}}\;\rule[-0.1em]{0.15em}{0.75em}\;\Varid{e}_{\mathrm{2}}\mskip1.5mu]\mskip1.5mu]} & \movesto & \ensuremath{\Conid{SX} [\mskip1.5mu CX [\mskip1.5mu \Varid{e}_{\mathrm{1}}\mskip1.5mu]\;\rule[-0.1em]{0.15em}{0.75em}\;CX [\mskip1.5mu \Varid{e}_{\mathrm{2}}\mskip1.5mu]\mskip1.5mu]} &  \\


\end{array} \\
\\[-0.3em]
\text{\textit{Note}: In the rules marked with a superscript $\alpha$, use $\alpha$-conversion to satisfy the side condition.}
\end{array}$
\iffullverse
$$\begin{array}{l}
  \textit{Effects} \\
    \begin{array}{lr@{\;}c@{\;}l}
      \rulename{succ-value} & \ensuremath{\textbf{succeeds}\;\Varid{v}} & \movesto & \ensuremath{\Varid{v}} \\
      \rulename{dec-fail} & \ensuremath{\textbf{decides}\;\textbf{fail}} & \movesto & \ensuremath{\textbf{fail}} \\
      \rulename{dec-value} & \ensuremath{\textbf{decides}\;\Varid{v}} & \movesto & \ensuremath{\Varid{v}} \\
    \end{array} \\
  \end{array}$$
\fi
\caption{\textbf{The \vcname: Rewrite rules}} \label{fig:rules}
\Description{A table that exhibits the rewrite rules for the Verse Calculus, groupesd into five categories: Application, Unification, Normalization, Garbage Collection, and Choice}
\end{figure}
}

\newcommand{\figureStructuralRules}{
  \begin{figure}
    $$
    \begin{array}{l}
      \textit{Structural rules} \\
      \begin{array}{lr@{\;}c@{\;}ll}
      \rulename{exi-swap} & \ensuremath{\exists\Varid{x}.\,\exists\Varid{y}.\,\Varid{e}} & \movesto & \ensuremath{\exists\Varid{y}.\,\exists\Varid{x}.\,\Varid{e}} \\
      \rulename{val-swap} & \ensuremath{eq;\,\Varid{x}\hspace{-0.14em}=\hspace{-0.14em}\Varid{v};\,\Varid{e}} & \movesto & \ensuremath{\Varid{x}\hspace{-0.14em}=\hspace{-0.14em}\Varid{v};\,eq;\,\Varid{e}} \\
      \text{\lennart{Plan D:}}\\
      \rulename{var-swap-subst} & \ensuremath{X [\mskip1.5mu \Varid{x}\hspace{-0.14em}=\hspace{-0.14em}\Varid{y}\mskip1.5mu]} & \movesto & \ensuremath{\subst{X}{\Varid{y}}{\Varid{x}} [\mskip1.5mu \Varid{y}\hspace{-0.14em}=\hspace{-0.14em}\Varid{x}\mskip1.5mu]} \\
      % \rulename{unify-swap} & |x == y)| & \movesto & |x == y| \\
      \end{array} \\

     \end{array}
     $$
    \caption{The Verse calculus: Structural Rules}
    \Description{The Verse calculus: Structural Rules}
      \label{fig:rules-structural}  \label{fig:structural}
   \end{figure}
}

